% ===========================================================================
% main.tex -- the English part of the thesis (pdfLaTeX + biber).
%
% The Hebrew summary/abstract live in hebrew_pages.tex (XeLaTeX). The two
% PDFs are merged by merge_thesis.py; run ./compile_thesis.sh to build all.
% See README.md for the whole picture and for the cross-document link rules.
% ===========================================================================
\documentclass[12pt,a4paper,oneside,onecolumn]{book}
\usepackage[utf8]{inputenc}
\usepackage[T1]{fontenc}
\usepackage{etoolbox}

% ---- Page layout (BGU guidelines: A4, page number top centre, no headers) ----
\usepackage[a4paper,twoside,inner=24mm,outer=24mm,top=28mm,bottom=32mm]{geometry}
% NOTE: `twoside` in geometry makes \cleardoublepage insert blank pages so that
% every chapter starts on an odd page (as in a printed, double-sided thesis).
% For a screen-only PDF add `openany` to the \documentclass options.
\usepackage{microtype}

\usepackage{fancyhdr}
\pagestyle{fancy}
\lhead{}\chead{\thepage}\rhead{}
\lfoot{}\cfoot{}\rfoot{}
\renewcommand{\headrulewidth}{0pt}
\fancypagestyle{plain}{\pagestyle{fancy}}
\setlength{\headheight}{15pt}

% ---- Maths and theorem environments ----
\usepackage{mathtools}
\usepackage{amssymb}
\usepackage[amsmath,thmmarks]{ntheorem}
\usepackage{thmtools}            % \begin{restatable}...: state in the intro, restate in a chapter

\newtheorem{theorem}{Theorem}
\newtheorem{lemma}[theorem]{Lemma}
\newtheorem{corollary}[theorem]{Corollary}
\newtheorem{proposition}[theorem]{Proposition}
\newtheorem{observation}[theorem]{Observation}
\newtheorem{definition}{Definition}

\theoremstyle{nonumberplain}
\theoremheaderfont{\itshape}
\theorembodyfont{\normalfont}
\theoremseparator{.}
\theoremsymbol{\ensuremath{\square}}
\newtheorem{proof}{Proof}
% back to defaults
\theoremstyle{plain}
\theoremheaderfont{\normalfont\bfseries}
\theorembodyfont{\itshape}
\theoremseparator{}
\theoremsymbol{}

\usepackage{chngcntr}
\counterwithin{theorem}{chapter}   % Theorem 2.1, Lemma 2.2, ...

\DeclareMathOperator*{\argmin}{arg\,min}
\newcommand{\NP}{\textnormal{NP}}
\newcommand{\PP}{\textnormal{P}}
\newcommand{\FPT}{\textnormal{FPT}}
\DeclareMathOperator*{\argmax}{arg\,max}

% ---- Misc ----
\usepackage[shortlabels]{enumitem}
\usepackage[ruled,vlined,algochapter]{algorithm2e}
\usepackage{booktabs}
\usepackage{array}
\usepackage{csquotes}
\usepackage{graphicx}
\graphicspath{{images/}}
\usepackage{tikz}
\usetikzlibrary{positioning,arrows.meta,calc}

% Code listings (the chapters show LaTeX snippets)
\usepackage{listings}
\lstset{basicstyle=\ttfamily\small,breaklines=true,columns=fullflexible,keepspaces=true,
        frame=single,framerule=0.3pt,xleftmargin=1em,aboveskip=0.8em,belowskip=0.8em,
        language=[LaTeX]TeX,texcsstyle=*\bfseries,commentstyle=\itshape\color{black!60},
        morekeywords={hypertarget,hyperlink,linkedsection,enref,hcite,en,englishTOCnumber,englishTOCpage}}
\setcounter{secnumdepth}{4}

% ---- Bibliography (biblatex + biber). The Hebrew document loads the same
%      .bib with the same style, so citation labels agree across both parts. ----
\usepackage[backend=biber,natbib=true,style=alphabetic,sorting=none,
            maxcitenames=2,maxbibnames=50,doi=true,url=true,isbn=false,eprint=false]{biblatex}
\renewcommand*{\multicitedelim}{\addcomma\space}
\addbibresource{bibliography/references.bib}
% Every bibliography entry also gets a named anchor equal to its bib key, so
% the Hebrew pages can \hyperlink{<bibkey>}{...} straight to an entry.
\AtEveryBibitem{\hypertarget{\thefield{entrykey}}{}}

% ---- Table of contents ----
\usepackage[nottoc]{tocbibind}
\usepackage{tocloft}
\renewcommand{\cftchapfont}{\bfseries}
\renewcommand{\cftchappagefont}{\bfseries}
\renewcommand{\cftchapdotsep}{\cftnodots}
\renewcommand{\cftsecfont}{\normalfont}
\renewcommand{\cftsecpagefont}{\normalfont}
\renewcommand{\cftsecdotsep}{1}
\setlength{\cftbeforechapskip}{14pt}
\setlength{\cftbeforesecskip}{0pt}

% ---- Links and cross-document references ----
% xr-hyper must come before hyperref. It lets this document read labels from
% the Hebrew build (prefix he-), which is how the English TOC learns the page
% numbers of the Hebrew summary and abstract.
\usepackage{xr-hyper}
\usepackage[urlcolor=blue!50!black,linkcolor=blue!50!black,citecolor=green!40!black,colorlinks]{hyperref}
\usepackage{lastpage}   % defines the LastPage label; hebrew_pages.tex continues numbering after it
\IfFileExists{build/hebrew_pages.aux}{\externaldocument[he-]{build/hebrew_pages}}{}

\usepackage[capitalise]{cleveref}

% \linkedsection{<key>}{<Title>}: a \section the Hebrew TOC can link to and read the
% number and page of. It expands to
%   \section[<Title>]{\hypertarget{target:sec:<key>}{}<Title>}\label{toc:sec:<key>}
% The anchor sits INSIDE the heading so a click lands on the heading itself even
% when the section starts a new page (an anchor placed before \section would end
% up at the bottom of the previous page). Add your own \label{sec:...} after it
% for \Cref if you like.
\newcommand{\linkedsection}[2]{\section[#2]{\hypertarget{target:sec:#1}{}#2}\label{toc:sec:#1}}
\Crefname{observation}{Observation}{Observations}

% ---- Spacing ----
\usepackage{setspace}
\setstretch{1.1}
\setlength{\parskip}{0.25em plus 0.1em minus 0.05em}
\setlength{\parindent}{2em}
\setlength{\emergencystretch}{3em}   % let TeX stretch lines rather than overflow them around long \texttt words

% ---- Who and what: title, author, supervisor, department, date ----
% ---------------------------------------------------------------------------
% thesis_info.tex -- the only file you must edit to make the thesis yours.
% It is read by BOTH documents (main.tex and hebrew_pages.tex), so the cover
% pages, the Hebrew cover and the abstracts all agree.
% ---------------------------------------------------------------------------

% Title
\newcommand{\thesistitle}{A Thesis Template for Ben-Gurion University}
\newcommand{\thesistitlehe}{תבנית לעבודת גמר באוניברסיטת בן-גוריון}

% Author
\newcommand{\thesisauthorname}{Israel Israeli}
\newcommand{\thesisauthornamehe}{ישראל ישראלי}

% Supervisor(s). For two supervisors, define \thesissecondsupervisername and
% \thesissecondsupervisernamehe as well and adjust Main_pages/cover_page.tex,
% Main_pages/title_page.tex and the Hebrew cover at the end of hebrew_pages.tex.
\newcommand{\thesissupervisername}{Prof. Max Musterman}
\newcommand{\thesissupervisernamehe}{פרופ' מקס מוסטרמן}

% Faculty and department
\newcommand{\thesisfaculty}{The Faculty of Engineering}
\newcommand{\thesisfacultyhe}{הפקולטה להנדסה}
\newcommand{\thesisdepartment}{The Department of Industrial Engineering and Management}
\newcommand{\thesisdepartmenthe}{המחלקה להנדסת תעשייה וניהול}

% Degree statement (PhD wording; change for an MSc thesis)
\newcommand{\thesisdegreetext}{Thesis submitted in partial fulfillment of the requirements\\for the degree of Doctor of Philosophy}
\newcommand{\thesisdegreetexthe}{חיבור זה מהווה חלק מהדרישות לקבלת התואר דוקטור לפילוסופיה}

% Submission date
\newcommand{\thesismonth}{September}
\newcommand{\thesismonthhe}{ספטמבר}
\newcommand{\thesisyear}{2026}


\begin{document}

% ---------------------------------------------------------------------------
% Cover and title page: no page numbers at all.
% ---------------------------------------------------------------------------
\pagenumbering{gobble}
\begingroup
  \makeatletter\let\ps@plain\ps@empty\makeatother
  \begin{titlepage}
    \begin{center}
        \vspace*{1cm}
        \includegraphics[width=0.1\textwidth]{Main_pages/bgu.png}\\
        {\Large\textbf{Ben-Gurion University of the Negev}}\\
        {\large \thesisfaculty}\\
        \thesisdepartment

        \vspace{2cm}

        {\Large \textbf{\thesistitle}}

        \vspace{1cm}

        \thesisdegreetext

        \vspace{1cm}

        \textbf{\thesisauthorname}

        \vspace{1cm}

        Under the supervision of \textbf{\thesissupervisername}

        \vfill

        \textbf{\thesismonth} \quad \textbf{\thesisyear}
    \end{center}
\end{titlepage}

  \begin{titlepage}
    \begin{center}
        \vspace*{1cm}
        \includegraphics[width=0.1\textwidth]{Main_pages/bgu.png}\\
        {\Large\textbf{Ben-Gurion University of the Negev}}\\
        {\large \thesisfaculty}\\
        \thesisdepartment

        \vspace{1cm}

        {\Large \textbf{\thesistitle}}

        \vspace{0.5cm}

        \thesisdegreetext

        \vspace{0.5cm}

        \textbf{\thesisauthorname}

        \vspace{0.5cm}

        Under the supervision of \textbf{\thesissupervisername}

        \vspace{4.5cm}

        \begin{flushleft}
        Signature of student: \( \rule{3cm}{0.15mm} \) \hfill Date: \( \rule{3cm}{0.15mm} \)\\
        \vspace{.65cm}
        Signature of supervisor: \( \rule{3cm}{0.15mm} \) \hfill Date: \( \rule{3cm}{0.15mm} \)\\
        \vspace{.25cm}
        Signature of chairperson of the\\committee for graduate studies: \( \rule{3cm}{0.15mm} \) \hfill Date: \( \rule{3cm}{0.15mm} \)
        \end{flushleft}

        \vfill

        \textbf{\thesismonth} \quad \textbf{\thesisyear}
    \end{center}
\end{titlepage}

\endgroup

% \frontmatter would insert a blank page after the title page in twoside
% mode; suppress that one time.
\makeatletter
\let\orig@cleardoublepage\cleardoublepage
\let\cleardoublepage\clearpage
\frontmatter
\let\cleardoublepage\orig@cleardoublepage
\makeatother
\pagenumbering{roman}

% ---------------------------------------------------------------------------
% Front matter: abstract, acknowledgements, table of contents.
% ---------------------------------------------------------------------------
% \noindent
{\Large\bfseries Research Student's Affidavit}

\vspace{1.2cm}

\noindent
I, \textbf{\thesisauthorname}, whose signature appears below, hereby declare that:

\vspace{0.5cm}

\begin{itemize}
    \setlength{\itemsep}{0.8em}
    \item I have written this thesis independently, except for the assistance and guidance provided by my thesis advisors.
    \item The scientific material included in this thesis is the product of my own research, conducted during the period of my studies as a research student.
\end{itemize}

\vspace{0.5cm}

\noindent
This thesis incorporates research materials produced in collaboration with others, excluding technical assistance commonly received in the course of experimental work. Accordingly, I am attaching an additional affidavit specifying my contribution and the contributions of the other participants in this research. This affidavit has been approved by all relevant parties and submitted with their consent.

\vfill

\noindent
\begin{tabular}{p{0.28\textwidth} p{0.32\textwidth} p{0.28\textwidth}}
\textbf{Date:} & \textbf{Student's Name:} & \textbf{Signature:} \\[1.5cm]
DD.MM.YYYY & \thesisauthorname & \hrulefill
\end{tabular}
   % if your department requires it
\phantomsection
\addcontentsline{toc}{chapter}{Abstract}
\thispagestyle{plain}
{\noindent\Large\textbf{Abstract}}

\vspace{3ex}

This document is a template for a thesis at Ben-Gurion University of the Negev, and at the same time its manual.
An English thesis and the Hebrew summary, abstract and cover that the university requires at the end of the bound volume are produced as one PDF whose links work in both directions: the English table of contents lists the Hebrew pages, the Hebrew table of contents shows the English chapter numbers and pages read from the English build, and Hebrew text can cite the bibliography and link to English theorems, figures and definitions while printing their numbers automatically.

The chapters explain the build process, the anchor and label conventions that connect the two parts, and the commands for writing Hebrew, and they exercise every construct a thesis needs along the way: numbered chapters and sections, a theorem stated in the introduction and restated with its proof, a definition, a figure, a table, an algorithm, code listings and citations.

Replace this abstract with your own.

\chapter*{Acknowledgements}

I thank my supervisor, \thesissupervisername, for guidance, patience, and the occasional well-timed question.
I thank the members of my committee, my colleagues at \thesisdepartment, and the people who read drafts of this thesis.

Financial support from the Kreitman School of Advanced Graduate Studies is gratefully acknowledged.

Finally, I thank my family and friends.


\setcounter{tocdepth}{2}
\tableofcontents
\clearpage
% \listoffigures \clearpage
% \listoftables  \clearpage

% ---------------------------------------------------------------------------
% Main matter. Every chapter follows the same four-line pattern:
%   \cleardoublepage \phantomsection \hypertarget{chapter:<key>}{} \chapter{..} \label{chapter:<key>}
% The \hypertarget is what the Hebrew TOC links to (\hyperlink{chapter:<key>}),
% the \label is what \Cref and the Hebrew TOC numbers/pages read.
% ---------------------------------------------------------------------------
\mainmatter

\cleardoublepage
\phantomsection
\hypertarget{chapter:intro}{}%
\chapter{Introduction}
\label{chapter:intro}
% ---------------------------------------------------------------------------
% content/introduction.tex
% Every chapter file is plain LaTeX. Sections use \linkedsection{<key>}{<Title>}
% (defined in main.tex) so that the Hebrew table of contents can link to them
% and print their number and page. See Chapter 3 for the rules.
% ---------------------------------------------------------------------------

\linkedsection{what}{What This Template Gives You}

A thesis at Ben-Gurion University is written in English and bound with a Hebrew summary, a Hebrew abstract and a Hebrew cover at its end, read from the back.
LaTeX cannot typeset both halves in one document without pain, so this template keeps them apart: \texttt{main.tex} is the English thesis, compiled with pdfLaTeX, and \texttt{hebrew\_pages.tex} is the Hebrew part, compiled with XeLaTeX.
A script compiles both and merges them into one PDF, the one you are reading, in which

\begin{itemize}
    \item the English table of contents lists the Hebrew pages with their real page numbers,
    \item the Hebrew table of contents shows the English chapter and section numbers and pages, read automatically from the English build, and each line of it is a link into the English part,
    \item Hebrew text can cite the bibliography and refer to any English theorem, figure or definition by a link that also prints its number,
    \item page numbers run continuously through both parts, and
    \item every link keeps working after the PDF has been opened and re-saved by any viewer.
\end{itemize}

This document is both the demonstration and the manual.
\Cref{chapter:build} explains the build, \Cref{chapter:links} the anchors that connect the two parts, and \Cref{chapter:hebrew} the commands for writing Hebrew.
Along the way, every construct a thesis needs appears at least once: a theorem stated here and restated with its proof later, a definition, a figure, a table, an algorithm, code listings, citations such as \cite{Karp72,GJ79} or \cite{downey2013fundamentals,cygan2015parameterized}, and cross-references with \verb|\Cref|.

The whole idea of the template fits in one statement, which \Cref{chapter:links} proves.

\begin{restatable}{theorem}{MainTheorem}
\label{thm:main}\hypertarget{target:thm:main}{}%
Every object in the English part that carries a named anchor can be linked to from the Hebrew pages, and the link survives any PDF viewer.
\end{restatable}

\linkedsection{quickstart}{Quick Start}

\begin{enumerate}
    \item Fill in \texttt{thesis\_info.tex}: title, author, supervisor, department, degree and date, in English and Hebrew. Both documents read this file, so the three cover pages and the abstracts agree.
    \item Write your chapters in \texttt{content/} and list them in \texttt{main.tex} by copying the chapter block shown in \Cref{sec:sections}.
    \item Write the Hebrew summary and abstract in \texttt{hebrew\_pages.tex}, and mirror your chapters in its table of contents (\Cref{sec:hebrewtoc}).
    \item Run \texttt{./compile\_thesis.sh} and open \texttt{build/thesis.pdf}. The merge step must print no warning (\Cref{sec:merge}).
\end{enumerate}

\Cref{tab:files} lists the files of the template.

\begin{table}[htb]
    \centering
    \caption{The files of the template.}
    \label{tab:files}
    \begin{tabular}{@{}l>{\raggedright\arraybackslash}p{0.58\textwidth}@{}}
        \toprule
        File & Role \\
        \midrule
        \texttt{thesis\_info.tex} & title, names, department, degree, date (English and Hebrew) \\
        \texttt{main.tex} & the English document: preamble, front matter, chapter list, bibliography \\
        \texttt{content/*.tex} & one file per chapter \\
        \texttt{Main\_pages/} & cover, signature page, abstract, acknowledgements, affidavit, logo \\
        \texttt{hebrew\_pages.tex} & the Hebrew document: Hebrew TOC, summary, abstract, Hebrew cover \\
        \texttt{bibliography/references.bib} & the bibliography, shared by both documents \\
        \texttt{compile\_thesis.sh} & builds everything into \texttt{build/thesis.pdf} \\
        \texttt{merge\_thesis.py} & concatenates the two PDFs and rewires the links \\
        \bottomrule
    \end{tabular}
\end{table}


\cleardoublepage
\phantomsection
\hypertarget{chapter:build}{}%
\chapter{The Build Process}
\label{chapter:build}
% ---------------------------------------------------------------------------
% content/build.tex -- how compile_thesis.sh works and why.
% ---------------------------------------------------------------------------

\linkedsection{script}{What the Script Does}

\texttt{compile\_thesis.sh} needs TeX Live (or MacTeX) with \texttt{latexmk}, \texttt{pdflatex}, \texttt{xelatex} and \texttt{biber}, and Python~3 with the \texttt{pikepdf} package.
It runs the three steps of \Cref{fig:pipeline}, all inside the git-ignored folder \texttt{build/}.

\begin{figure}[htb]
\hypertarget{target:fig:pipeline}{}%
    \centering
    \resizebox{\textwidth}{!}{%
    \begin{tikzpicture}[
        box/.style={draw,rounded corners,minimum height=9mm,minimum width=27mm,align=center,font=\small},
        arr/.style={-{Stealth[length=2.5mm]},thick}]
        \node[box] (en) {\texttt{main.tex}\\pdfLaTeX + biber};
        \node[box,below=9mm of en] (he) {\texttt{hebrew\_pages.tex}\\XeLaTeX + biber};
        \node[box,right=22mm of en] (enpdf) {\texttt{build/main.pdf}};
        \node[box,right=22mm of he] (hepdf) {\texttt{build/hebrew\_pages.pdf}};
        \node[box,right=22mm of $(enpdf)!0.5!(hepdf)$] (merge) {\texttt{merge\_thesis.py}};
        \node[box,right=12mm of merge] (out) {\texttt{build/thesis.pdf}};
        \draw[arr] (en) -- (enpdf);
        \draw[arr] (he) -- (hepdf);
        \draw[arr] (enpdf) -- (merge);
        \draw[arr] (hepdf) -- (merge);
        \draw[arr] (merge) -- (out);
        \draw[arr,dashed] (en.south) to[bend right=25] node[left,font=\scriptsize,align=right] {page numbers,\\\texttt{LastPage}} (he.north);
        \draw[arr,dashed] (he.north east) to[bend right=25] node[right,font=\scriptsize,align=left] {Hebrew\\page numbers} (en.south east);
    \end{tikzpicture}}
    \caption{The build. Solid arrows are the three steps of the script; dashed arrows are the labels each document reads from the other's auxiliary file, which is why the first two steps are repeated until both settle.}
    \label{fig:pipeline}
\end{figure}

\begin{enumerate}
    \item \texttt{latexmk -pdf main.tex} produces \texttt{build/main.pdf}. \texttt{latexmk} runs biber and repeats pdfLaTeX until the references are stable.
    \item \texttt{latexmk -xelatex hebrew\_pages.tex} produces \texttt{build/hebrew\_pages.pdf}.
    \item \texttt{merge\_thesis.py} appends the Hebrew pages to the English PDF and rewires the cross-document links (\Cref{sec:merge}). The result is \texttt{build/thesis.pdf}.
\end{enumerate}

\texttt{./compile\_thesis.sh clean} deletes \texttt{build/}. A clean build takes a few minutes; later builds only redo what changed.

\linkedsection{settle}{Why the First Two Steps Repeat}

Each document reads labels from the other.
The Hebrew document continues the page numbering after the English \texttt{LastPage} label and prints the English chapter numbers and pages in its table of contents.
The English document, in turn, needs the page numbers of the Hebrew summary and abstract for its own table of contents.
Both documents load the \texttt{xr-hyper} package, which imports the labels of another document's \texttt{.aux} file under a prefix: \texttt{en-} for English labels read by the Hebrew document, \texttt{he-} for Hebrew labels read by the English one.

Because each side depends on the other, the script compiles English, then Hebrew, then checks whether either \texttt{.aux} file changed, and repeats the pair until neither does.
On a clean build the first English pass prints page~0 for the Hebrew entries, since no Hebrew \texttt{.aux} exists yet; the second pass fixes it.
From the second pass on, the script passes \texttt{-g} to \texttt{latexmk}, because \texttt{latexmk} does not know that the other document's \texttt{.aux} file is a dependency and would otherwise report that there is nothing to do.

\linkedsection{troubleshooting}{Fonts, Overleaf and Troubleshooting}

\paragraph{Fonts.} The Hebrew pages use \emph{Linux Libertine O} for Latin text, \emph{David CLM} for Hebrew and \emph{Latin Modern Roman} for Latin fragments inside Hebrew paragraphs. All three ship with TeX Live (packages \texttt{libertine}, \texttt{culmus}, \texttt{lm}), but XeTeX on macOS only finds a font by name if macOS itself knows it. If the build stops with \texttt{The font "Linux Libertine O" cannot be found}, run \texttt{./compile\_thesis.sh install-fonts} once; it copies the font files into \texttt{\textasciitilde/Library/Fonts}. On Linux and on Overleaf this is not needed.

\paragraph{Overleaf.} Overleaf cannot run the merge script. Compile \texttt{main.tex} with pdfLaTeX and \texttt{hebrew\_pages.tex} with XeLaTeX (two projects, or one project with the main file switched), download both PDFs and merge locally with \texttt{python3 merge\_thesis.py english.pdf hebrew.pdf out.pdf}.

\paragraph{Symptoms and causes.}
\begin{itemize}
    \item \emph{Hebrew Summary~0 in the English TOC}: the English part was not rebuilt after the Hebrew part. Run the script again rather than \texttt{latexmk} by hand.
    \item \emph{Merge warning about names that exist in neither PDF}: a Hebrew \verb|\hyperlink| has no matching anchor in the English source. See \Cref{sec:merge}.
    \item \emph{Nonsense errors such as \texttt{Missing \textbackslash begin\{document\}} in \texttt{build/main.aux}}: a build was interrupted or two ran at once and an \texttt{.aux} file is truncated. Run \texttt{./compile\_thesis.sh clean} and build again.
    \item \emph{biber: extracting arm64 binary with lipo failed} (macOS, after a system update): the universal \texttt{biber} shipped with MacTeX unpacks itself with \texttt{lipo}, which newer command line tools reject. Extract the slice once by hand, \texttt{lipo -thin arm64 \$(which biber) -output \textasciitilde/bin/biber}, and put \texttt{\textasciitilde/bin} first on your \texttt{PATH}.
    \item \emph{Package bidi error about xcolor}: in \texttt{hebrew\_pages.tex}, any package that loads \texttt{xcolor} must come before \texttt{bidi} and \texttt{polyglossia}. The template loads \texttt{xcolor} first for this reason.
\end{itemize}


\cleardoublepage
\phantomsection
\hypertarget{chapter:links}{}%
\chapter{Anchors and Links Between the Two Parts}
\label{chapter:links}
% ---------------------------------------------------------------------------
% content/links.tex -- the anchor and label conventions.
% ---------------------------------------------------------------------------

\linkedsection{anchors}{Anchors and Labels}

Two different things connect the Hebrew pages to the English ones, and it helps to keep them apart.

\begin{definition}
\label{def:anchor}\hypertarget{target:def:anchor}{}%
An \emph{anchor} is a named position in a PDF, created by \verb|\hypertarget{name}{}| and reached by \verb|\hyperlink{name}{text}|. A \emph{label} is a LaTeX name, created by \verb|\label{name}|, from which \verb|\ref| and \verb|\pageref| print a number and a page.
\end{definition}

Anchors are what a click follows; labels are what gets printed.
The Hebrew document can use both across the document boundary: labels through \texttt{xr-hyper}, and anchors through the merge step, since an anchor's name is stored in the PDF.
The template's convention is to give every object the Hebrew pages may refer to both an anchor and a label, named as in \Cref{tab:naming}.

\begin{table}[htb]
    \centering
    \caption{Naming conventions. The Hebrew document reads labels with the prefix \texttt{en-}.}
    \label{tab:naming}\hypertarget{target:tab:naming}{}%
    \begin{tabular}{@{}>{\raggedright\arraybackslash}p{0.2\textwidth}>{\raggedright\arraybackslash}p{0.4\textwidth}>{\raggedright\arraybackslash}p{0.3\textwidth}@{}}
        \toprule
        Object & Anchor (for \verb|\hyperlink|) & Label (for \verb|\ref*|, \verb|\pageref*|) \\
        \midrule
        Chapter & \texttt{chapter:<key>} & \texttt{chapter:<key>} \\
        Section & \texttt{target:sec:<key>} & \texttt{toc:sec:<key>} \\
        Theorem, figure, table, definition & \texttt{target:thm:<key>}, \texttt{target:fig:<key>}, \ldots & \texttt{thm:<key>}, \texttt{fig:<key>}, \ldots \\
        Bibliography & \texttt{chapter:bib} & \texttt{toc:bibliography} \\
        Bibliography entry & its bib key & \\
        Hebrew summary, abstract & \texttt{chapter:hebrew:summary}, \texttt{chapter:hebrew:abstract} & \texttt{toc:hebrew:summary}, \texttt{toc:hebrew:abstract} \\
        \bottomrule
    \end{tabular}
\end{table}

\linkedsection{sections}{Chapters and Sections}
\label{sec:sections}

Every chapter in \texttt{main.tex} follows the same block.
The \verb|\phantomsection| gives hyperref a correct bookmark, the \verb|\hypertarget| is the anchor the Hebrew TOC links to, and the \verb|\label| is where it reads the chapter number and page.

\begin{lstlisting}
\cleardoublepage
\phantomsection
\hypertarget{chapter:build}{}%
\chapter{The Build Process}
\label{chapter:build}
% ---------------------------------------------------------------------------
% content/build.tex -- how compile_thesis.sh works and why.
% ---------------------------------------------------------------------------

\linkedsection{script}{What the Script Does}

\texttt{compile\_thesis.sh} needs TeX Live (or MacTeX) with \texttt{latexmk}, \texttt{pdflatex}, \texttt{xelatex} and \texttt{biber}, and Python~3 with the \texttt{pikepdf} package.
It runs the three steps of \Cref{fig:pipeline}, all inside the git-ignored folder \texttt{build/}.

\begin{figure}[htb]
\hypertarget{target:fig:pipeline}{}%
    \centering
    \resizebox{\textwidth}{!}{%
    \begin{tikzpicture}[
        box/.style={draw,rounded corners,minimum height=9mm,minimum width=27mm,align=center,font=\small},
        arr/.style={-{Stealth[length=2.5mm]},thick}]
        \node[box] (en) {\texttt{main.tex}\\pdfLaTeX + biber};
        \node[box,below=9mm of en] (he) {\texttt{hebrew\_pages.tex}\\XeLaTeX + biber};
        \node[box,right=22mm of en] (enpdf) {\texttt{build/main.pdf}};
        \node[box,right=22mm of he] (hepdf) {\texttt{build/hebrew\_pages.pdf}};
        \node[box,right=22mm of $(enpdf)!0.5!(hepdf)$] (merge) {\texttt{merge\_thesis.py}};
        \node[box,right=12mm of merge] (out) {\texttt{build/thesis.pdf}};
        \draw[arr] (en) -- (enpdf);
        \draw[arr] (he) -- (hepdf);
        \draw[arr] (enpdf) -- (merge);
        \draw[arr] (hepdf) -- (merge);
        \draw[arr] (merge) -- (out);
        \draw[arr,dashed] (en.south) to[bend right=25] node[left,font=\scriptsize,align=right] {page numbers,\\\texttt{LastPage}} (he.north);
        \draw[arr,dashed] (he.north east) to[bend right=25] node[right,font=\scriptsize,align=left] {Hebrew\\page numbers} (en.south east);
    \end{tikzpicture}}
    \caption{The build. Solid arrows are the three steps of the script; dashed arrows are the labels each document reads from the other's auxiliary file, which is why the first two steps are repeated until both settle.}
    \label{fig:pipeline}
\end{figure}

\begin{enumerate}
    \item \texttt{latexmk -pdf main.tex} produces \texttt{build/main.pdf}. \texttt{latexmk} runs biber and repeats pdfLaTeX until the references are stable.
    \item \texttt{latexmk -xelatex hebrew\_pages.tex} produces \texttt{build/hebrew\_pages.pdf}.
    \item \texttt{merge\_thesis.py} appends the Hebrew pages to the English PDF and rewires the cross-document links (\Cref{sec:merge}). The result is \texttt{build/thesis.pdf}.
\end{enumerate}

\texttt{./compile\_thesis.sh clean} deletes \texttt{build/}. A clean build takes a few minutes; later builds only redo what changed.

\linkedsection{settle}{Why the First Two Steps Repeat}

Each document reads labels from the other.
The Hebrew document continues the page numbering after the English \texttt{LastPage} label and prints the English chapter numbers and pages in its table of contents.
The English document, in turn, needs the page numbers of the Hebrew summary and abstract for its own table of contents.
Both documents load the \texttt{xr-hyper} package, which imports the labels of another document's \texttt{.aux} file under a prefix: \texttt{en-} for English labels read by the Hebrew document, \texttt{he-} for Hebrew labels read by the English one.

Because each side depends on the other, the script compiles English, then Hebrew, then checks whether either \texttt{.aux} file changed, and repeats the pair until neither does.
On a clean build the first English pass prints page~0 for the Hebrew entries, since no Hebrew \texttt{.aux} exists yet; the second pass fixes it.
From the second pass on, the script passes \texttt{-g} to \texttt{latexmk}, because \texttt{latexmk} does not know that the other document's \texttt{.aux} file is a dependency and would otherwise report that there is nothing to do.

\linkedsection{troubleshooting}{Fonts, Overleaf and Troubleshooting}

\paragraph{Fonts.} The Hebrew pages use \emph{Linux Libertine O} for Latin text, \emph{David CLM} for Hebrew and \emph{Latin Modern Roman} for Latin fragments inside Hebrew paragraphs. All three ship with TeX Live (packages \texttt{libertine}, \texttt{culmus}, \texttt{lm}), but XeTeX on macOS only finds a font by name if macOS itself knows it. If the build stops with \texttt{The font "Linux Libertine O" cannot be found}, run \texttt{./compile\_thesis.sh install-fonts} once; it copies the font files into \texttt{\textasciitilde/Library/Fonts}. On Linux and on Overleaf this is not needed.

\paragraph{Overleaf.} Overleaf cannot run the merge script. Compile \texttt{main.tex} with pdfLaTeX and \texttt{hebrew\_pages.tex} with XeLaTeX (two projects, or one project with the main file switched), download both PDFs and merge locally with \texttt{python3 merge\_thesis.py english.pdf hebrew.pdf out.pdf}.

\paragraph{Symptoms and causes.}
\begin{itemize}
    \item \emph{Hebrew Summary~0 in the English TOC}: the English part was not rebuilt after the Hebrew part. Run the script again rather than \texttt{latexmk} by hand.
    \item \emph{Merge warning about names that exist in neither PDF}: a Hebrew \verb|\hyperlink| has no matching anchor in the English source. See \Cref{sec:merge}.
    \item \emph{Nonsense errors such as \texttt{Missing \textbackslash begin\{document\}} in \texttt{build/main.aux}}: a build was interrupted or two ran at once and an \texttt{.aux} file is truncated. Run \texttt{./compile\_thesis.sh clean} and build again.
    \item \emph{biber: extracting arm64 binary with lipo failed} (macOS, after a system update): the universal \texttt{biber} shipped with MacTeX unpacks itself with \texttt{lipo}, which newer command line tools reject. Extract the slice once by hand, \texttt{lipo -thin arm64 \$(which biber) -output \textasciitilde/bin/biber}, and put \texttt{\textasciitilde/bin} first on your \texttt{PATH}.
    \item \emph{Package bidi error about xcolor}: in \texttt{hebrew\_pages.tex}, any package that loads \texttt{xcolor} must come before \texttt{bidi} and \texttt{polyglossia}. The template loads \texttt{xcolor} first for this reason.
\end{itemize}

\end{lstlisting}

Sections use the helper \verb|\linkedsection|, defined in \texttt{main.tex}:

\begin{lstlisting}
\linkedsection{script}{What the Script Does}
% expands to
\section[What the Script Does]{\hypertarget{target:sec:script}{}What the Script Does}
\label{toc:sec:script}
\end{lstlisting}

The anchor is placed \emph{inside} the heading on purpose. An anchor written on the line before \verb|\section| is set at the end of the preceding paragraph, so whenever the section happens to start a new page, a click lands at the bottom of the previous one.
Add your own \verb|\label{sec:...}| after \verb|\linkedsection| for \verb|\Cref|, as this section does.

An unnumbered chapter, such as the conclusion of this document, has no number to read, so it needs only the anchor and a label, plus \verb|\addcontentsline| to appear in the English TOC.
\Cref{alg:chapter} summarises the steps for adding a chapter that both tables of contents can see.

\begin{algorithm}[htb]
\caption{Adding a chapter}
\label{alg:chapter}
\KwIn{a chapter title, a short key, a Hebrew title}
write the chapter in \texttt{content/<key>.tex}, using \texttt{\textbackslash linkedsection} for its sections\;
copy the chapter block into \texttt{main.tex} with \texttt{chapter:<key>} as anchor and label\;
\ForEach{section in the chapter}{
    add a line to the Hebrew TOC with \texttt{toc:sec:<key>} and \texttt{target:sec:<key>} (\Cref{sec:hebrewtoc})\;
}
add a line for the chapter itself with \texttt{chapter:<key>}\;
run \texttt{./compile\_thesis.sh} and check that the merge step prints no warning\;
\end{algorithm}

\linkedsection{objects}{Theorems, Figures, Tables and Definitions}
\label{sec:objects}

Anything the Hebrew text should be able to point at gets an anchor right after its label.
For a theorem, the anchor goes at the start of the body; for a figure or table, right after \verb|\begin{figure}| or next to the \verb|\label|.

\begin{lstlisting}
\begin{restatable}{theorem}{MainTheorem}
\label{thm:main}\hypertarget{target:thm:main}{}%
Every object in the English part that carries ...
\end{restatable}

\begin{figure}[htb]
\hypertarget{target:fig:pipeline}{}%
    ...
    \label{fig:pipeline}
\end{figure}
\end{lstlisting}

The theorem above uses \texttt{restatable} from the \texttt{thmtools} package: it is stated in the introduction and restated here with \verb|\MainTheorem*|, keeping its number, so the two statements cannot drift apart.

\MainTheorem*

\begin{proof}
Let $x$ be an object with anchor \texttt{target:x}.
The Hebrew source writes \verb|\hyperlink{target:x}{...}|, which XeLaTeX turns into a link whose destination is the \emph{name} \texttt{target:x}.
Inside the Hebrew PDF alone this name is undefined.
After the merge, the combined PDF contains both the link and, from the English half, the anchor of that name (\Cref{sec:merge}).
The merge script then replaces the name in the link by the page and position of the anchor, so the link no longer depends on the name table, which some viewers discard when they re-save a file.
\end{proof}

\linkedsection{merge}{The Merge Step}
\label{sec:merge}

\texttt{merge\_thesis.py} appends the pages of \texttt{build/hebrew\_pages.pdf} to \texttt{build/main.pdf}, collects the named anchors of both files, and rewrites every link that refers to a name into an explicit destination, as in the proof above.
Names defined in both files are taken from the English half, except the \texttt{chapter:hebrew:*} anchors, which are taken from the Hebrew half.

The script prints how many links it resolved and, if any Hebrew link uses a name that neither document defines, a warning listing those names.
A clean thesis prints no warning.
The fix is to add the missing \verb|\hypertarget| in the English source.
As a last resort, the \texttt{ALIASES} table at the top of the script can map a name to one of hyperref's automatic anchors, such as \texttt{section.2.1} or \texttt{theorem.3.1}; those names change whenever numbering changes, so prefer a real anchor.

Citations need no anchors of their own: both documents load the same \texttt{.bib} file with the same \texttt{biblatex} options, and after the merge a Hebrew citation link resolves to the entry in the English bibliography (\Cref{sec:hcite}).
In addition, \texttt{main.tex} gives every bibliography entry an anchor named by its bib key, so \verb|\hyperlink{Karp72}{...}| works from either side.


\cleardoublepage
\phantomsection
\hypertarget{chapter:hebrew}{}%
\chapter{Writing the Hebrew Pages}
\label{chapter:hebrew}
% ---------------------------------------------------------------------------
% content/hebrewpages.tex -- the commands available in hebrew_pages.tex.
% ---------------------------------------------------------------------------

\linkedsection{rtl}{Right-to-Left Text and Latin Fragments}

\texttt{hebrew\_pages.tex} sets Hebrew as its default language with the \texttt{polyglossia} package, so ordinary paragraphs are typeset right to left and headings are aligned to the right.
Inline mathematics such as \verb|$f(k)\cdot n^{O(1)}$| works unchanged.
A Latin word or phrase inside a Hebrew sentence, however, must be marked, or the bidirectional algorithm will reorder it. The template defines

\begin{lstlisting}
\newcommand{\en}[1]{\beginL{\fontspec{Latin Modern Roman}#1}\endL}
\end{lstlisting}

so that \verb|\en{NP}| or \verb|\en{(fixed-parameter tractable)}| is typeset left to right, in a Latin font.
The year on the Hebrew cover page is set the same way, with \verb|\en{\thesisyear}|.

\linkedsection{hcite}{Citations: the \texttt{\textbackslash hcite} Command}
\label{sec:hcite}

A citation inside Hebrew text is a Latin fragment too, with one more requirement: it must be a link into the English bibliography.
The template's command is

\begin{lstlisting}
\newcommand{\hcite}[1]{\mbox{\LR{{\fontspec{Latin Modern Roman}\cite{#1}}}}}
\end{lstlisting}

and it is used exactly like \verb|\cite|: \verb|\hcite{Karp72,GJ79}|.
Three things make it work.

\begin{itemize}
    \item \verb|\LR| (from the \texttt{bidi} package) and the Latin font keep the citation label readable and left to right; \verb|\mbox| keeps it on one line.
    \item \texttt{hebrew\_pages.tex} loads the same \texttt{.bib} file with the same \texttt{biblatex} options as \texttt{main.tex}. The \texttt{alphabetic} style derives its labels from author names and years, so the Hebrew document prints \texttt{[Kar72]} for the same entry the English bibliography lists as \texttt{[Kar72]}, even though the Hebrew document prints no bibliography of its own.
    \item \texttt{biblatex} makes each citation a link to the anchor \texttt{cite.0@<key>}. In the Hebrew PDF alone that anchor is missing; after the merge it is found in the English bibliography and the link is rewired to it (\Cref{sec:merge}).
\end{itemize}

One rule follows from this: an entry can only be linked to if it appears in the English bibliography, that is, if it is cited (or listed with \verb|\nocite|) somewhere in \texttt{main.tex}. An entry cited only from Hebrew still prints its label, but the merge step reports its link as unresolved.

If you switch the English bibliography to a numeric style, switch the Hebrew document too, and keep in mind that numeric labels then depend on citation order and will differ between the two documents; the links still work, the printed numbers will not match. The alphabetic style avoids the problem.

\linkedsection{enref}{References to English Results: the \texttt{\textbackslash enref} Command}

To refer to an English theorem, figure, definition or table from Hebrew text, use

\begin{lstlisting}
\newcommand{\enref}[2]{\hyperlink{target:#1}{\englishTOCnumber{#2}}}
% usage, inside Hebrew text:  (<Hebrew word for Theorem> \enref{thm:main}{thm:main})
\end{lstlisting}

The first argument is the anchor name without its \texttt{target:} prefix, the second is the label.
The command prints the object's number, read from the English build, and makes it a link to the anchor, so the Hebrew text needs no hard-coded numbers and does not break when the English part is renumbered.
The Hebrew summary at the end of this document refers to \Cref{thm:main}, \Cref{def:anchor}, \Cref{fig:pipeline} and \Cref{tab:naming} this way.

Behind \verb|\enref| are the two macros that also drive the Hebrew table of contents:

\begin{lstlisting}
\newcommand{\englishTOCpage}[1]{\LR{\pageref*{en-#1}}}
\newcommand{\englishTOCnumber}[1]{\LR{\ref*{en-#1}}}
\end{lstlisting}

They print a page or a number from an \texttt{en-} label, left to right, and deliberately without a link of their own (the starred forms): the enclosing \verb|\hyperlink| is the link, and the merge script expects exactly one destination per link.

\linkedsection{hebrewtoc}{The Hebrew Table of Contents}
\label{sec:hebrewtoc}

The Hebrew TOC is written by hand, one line per English chapter or section, in \texttt{hebrew\_pages.tex}.
Each line combines the three ingredients above: a number, a linked Hebrew title and a page.

\begin{lstlisting}
\englishTOCnumber{chapter:build}.\ \hyperlink{chapter:build}{<Hebrew chapter title>}
\dotfill \englishTOCpage{chapter:build}\\[0.2em]
\hspace{2em}\englishTOCnumber{toc:sec:script}\ \hyperlink{target:sec:script}{<Hebrew section title>}
\dotfill \englishTOCpage{toc:sec:script}\\[0.1em]
\end{lstlisting}

Unnumbered chapters (the conclusion, the bibliography) omit the number.
The last two lines, for the Hebrew summary and abstract, use local labels, since those pages are in the Hebrew document itself; the English TOC reads the same labels through the \texttt{he-} prefix to list them on its side.
The order of the lines, and the Hebrew titles, are yours to maintain; the numbers and pages are not.


% An unnumbered chapter: it needs \addcontentsline to appear in the TOC.
\cleardoublepage
\hypertarget{chapter:conclusion}{}%
\chapter*{Conclusion}
\addcontentsline{toc}{chapter}{Conclusion}
\label{chapter:conclusion}
% ---------------------------------------------------------------------------
% content/conclusion.tex -- an unnumbered chapter (see main.tex for how it is
% added to the table of contents and anchored for the Hebrew TOC).
% ---------------------------------------------------------------------------

This template takes care of the mechanics that make a BGU thesis awkward to typeset: an English body and a Hebrew back part in two documents, compiled by two engines, joined into a single PDF whose tables of contents, page numbers and links agree.
The conventions are few.
Chapters are listed in \texttt{main.tex} with an anchor and a label, sections are written with \verb|\linkedsection|, and any theorem, figure or definition the Hebrew text refers to gets an anchor next to its label.
On the Hebrew side, \verb|\en| marks Latin fragments, \verb|\hcite| cites, and \verb|\enref| links to an English result and prints its number.
The script compiles both parts until they agree and merges them, and its merge step tells you if a link has nowhere to go.

To turn the template into your thesis, fill in \texttt{thesis\_info.tex}, replace the chapters in \texttt{content/} and the abstract and acknowledgements in \texttt{Main\_pages/}, write the Hebrew summary and abstract in \texttt{hebrew\_pages.tex}, and keep its table of contents in step with your chapters.
Include the affidavit if your department requires it, and decide whether the thesis will be printed double-sided or read on screen, which determines whether chapters should open on odd pages.
Before submission, build from a clean \texttt{build/} folder, make sure the merge step prints no warning, and open the PDF in the viewer your readers will use to click a few links in both directions.


% ---------------------------------------------------------------------------
% Back matter: bibliography, then the TOC lines for the Hebrew pages.
% ---------------------------------------------------------------------------
\backmatter

\phantomsection
\addcontentsline{toc}{chapter}{Bibliography}
\hypertarget{chapter:bib}{}
\AtNextBibliography{\label{toc:bibliography}}   % page of the bibliography heading, read by the Hebrew TOC
\printbibliography

% The Hebrew pages are appended after this document by merge_thesis.py, so
% their TOC entries are written by hand. The page numbers come from labels in
% hebrew_pages.tex (via xr-hyper, prefix he-); the last argument is the
% hyperref anchor the entry links to, defined in the Hebrew document.
\addtocontents{toc}{\protect\contentsline{chapter}{Hebrew Summary}{\getpagerefnumber{he-toc:hebrew:summary}}{chapter:hebrew:summary}}
\addtocontents{toc}{\protect\contentsline{chapter}{Hebrew Abstract}{\getpagerefnumber{he-toc:hebrew:abstract}}{chapter:hebrew:abstract}}

\clearpage
\end{document}
